\section{Green Agent：把 Benchmark 实现成可运行评估服务}

前面的 benchmark 都包含 task、environment、participant 和 evaluator。课程项目将 evaluator/orchestrator 封装为 Green Agent，把被测系统称为 White Agent。颜色只是角色名：Green Agent 代表评估方，负责准备测试、驱动交互、验证 outcome 和提交报告；White Agent 是 participant，不应获得 hidden test、ground truth 或 evaluator 内部状态。

\subsection{Testing my Agent：从静态题目到多 Agent 协议}

slide 65 用考试类比说明，评估者必须独立于考生。Green Agent 不仅发送问题，还可能启动环境、提供工具接口、控制预算、收集消息并在结束后评分。White Agent 则实现待测能力。平台还可能有 participant agents 或环境服务。角色分离是防止答案泄漏和保证可替换性的基础。

\lecturefigure{slides-images/slide-065.png}{官方 slide 65：Green evaluator、White participant 与测试角色}

\begin{importantbox}{Green/White 信任边界}
White Agent 只能看到完成任务所需的公开接口；Green Agent 可访问 hidden cases、ground truth 和评分逻辑，但不能修改 participant 的输出来“帮助通过”。所有消息、工具调用和环境变更应记录。评估代码与 participant code 分仓或分进程运行，可降低意外泄漏和依赖冲突。
\end{importantbox}

\begin{warningbox}{Evaluator 也可能成为攻击面}
White Agent 的输出可能包含 prompt injection、超长内容或恶意参数，诱导 Green Agent 泄漏答案或执行危险工具。Green Agent 应把 participant 内容当不可信输入，使用结构化协议、最小权限、超时/资源限制和独立 verifier。评估者不是天然安全的“上帝进程”。
\end{warningbox}

\subsection{职责与完整 assessment flow}

slides 67--69 逐步展开 Green Agent 职责：准备 environment、分发任务、驱动/观察 participant、验证结果并把报告提交平台。完整 flow 从平台确认 agents ready 开始，Green Agent 再 orchestration interaction，最终 metrics 被计算并上传。这个流程把 benchmark 论文中的一次实验变成可重复服务。

\lecturefigure{slides-images/slide-067.png}{官方 slide 67：Green Agent 的准备、交互、环境、验证与报告职责}
\lecturefigure{slides-images/slide-068.png}{官方 slide 68：Green Agent 作为 evaluator 和 orchestrator}
\lecturefigure{slides-images/slide-069.png}{官方 slide 69：平台、Green、White、environment 的完整 assessment flow}

\begin{knowledgebox}{读图：控制平面与数据平面}
平台负责 agent discovery、任务生命周期和结果登记，可视为控制平面；Green/White 消息、工具调用和环境状态是数据平面。控制平面不应理解任务语义，Green Agent 不应承担平台账户和全局调度。边界清楚后，同一 Green benchmark 才能测试不同 White Agent，同一 White Agent 也能接入多个 benchmark。
\end{knowledgebox}

\begin{importantbox}{一次评估 run 的状态机}
建议使用状态：CREATED $\rightarrow$ READY $\rightarrow$ RUNNING $\rightarrow$ VERIFYING $\rightarrow$ REPORTED，并为 TIMEOUT、FAILED、CANCELLED 定义终态。每次状态转换保存时间、agent 版本、environment digest 和错误原因。没有显式状态机，网络重试容易重复启动任务或重复计分。
\end{importantbox}

\subsection{Green Agent 如何构建}

slides 70--71 给出典型组成：task loader/selector、orchestrator、environment、metric modules 与 report writer。prompt-based toolkit 可以帮助开发者描述任务和角色，但真实 benchmark 仍需代码初始化环境、处理协议和执行 verifier。课程提醒，有些项目需要 2--3 个 Agent，有些还包含 user simulator、service simulator 或额外 judge。

\lecturefigure{slides-images/slide-070.png}{官方 slide 70：Green Agent 的任务、orchestrator、环境与 metric 组件}
\lecturefigure{slides-images/slide-071.png}{官方 slide 71：Prompt-based toolkit 与多 Agent 拓扑}

\begin{knowledgebox}{组件职责表}
\begin{longtable}{p{3cm}p{7.7cm}}
\toprule
\textbf{组件} & \textbf{职责} \\
\midrule
Task loader & 读取公开/hidden case，生成 run-specific task，不泄漏答案 \\
Environment manager & 创建、重置、销毁隔离环境，固定依赖与 seed \\
Orchestrator & 发送 kickoff、路由消息、执行预算/超时、判定结束 \\
Verifier/metric & 从 outcome/trajectory 计算指标并返回可解释证据 \\
Reporter & 生成结构化结果、版本元数据和失败日志，提交平台 \\
\bottomrule
\end{longtable}
\end{knowledgebox}

\textbf{课堂提示：}Green Agent 的智能程度取决于 benchmark。可执行 unit test 足够时，evaluator 应尽量确定；开放任务可能需要 judge Agent。不要为了“Agentic”而让所有组件都由 LLM 决策，确定性代码更便宜、可复现，也更不易被 participant 操纵。

\subsection{本章小结}

Green Agent 是 benchmark 的运行时封装：它管理环境与交互、保护 hidden evidence、计算 metrics 并报告。角色、权限、状态机和组件边界决定评估能否可信复现。

